% Part: many-valued-logic
% Chapter: three-valued-logics
% Section: lukasiewicz

\documentclass[../../../include/open-logic-section]{subfiles}

\begin{document}

\olfileid[tr]{mvl}{thr}{luk}

\olsection{\L ukasiewicz Mantığı}

Çok değerli mantığa ilişkin yayımlanmış ve ayrıntılı biçimde geliştirilmiş
ilk önerilerden biri, 1921'de Polonyalı filozof Jan \L ukasiewicz tarafından
ortaya konmuştur. \L ukasiewicz'i Aristoteles'in deniz savaşı problemi
güdülemişti: \emph{Bugün}, “Yarın bir deniz savaşı olacak” tümcesi ne doğru
ne de yanlış görünür; doğruluk değeri henüz belirlenmemiştir. \L ukasiewicz
bu tür “geleceğe ilişkin olumsal” tümceler için üçüncü bir doğruluk değeri
getirmeyi önerdi.
\begin{quote}
Önümüzdeki yılın belirli bir anında, örneğin 21 Aralık günü öğleyin
Varşova'da bulunmamın şu anda ne olumlu ne de olumsuz biçimde belirlenmiş
olduğunu çelişkiye düşmeden varsayabilirim. Dolayısıyla o anda Varşova'da
bulunmam mümkündür, fakat zorunlu değildir. Bu varsayım altında “Önümüzdeki
yıl 21 Aralık günü öğleyin Varşova'da olacağım” önermesi şu anda ne doğru ne
de yanlış olabilir. Çünkü şimdi doğru olsaydı gelecekte Varşova'da bulunmam
zorunlu olurdu ki bu varsayımla çelişir. Öte yandan şimdi yanlış olsaydı
gelecekte Varşova'da bulunmam olanaksız olurdu ki bu da varsayımla çelişir.
Bu nedenle ele alınan önerme şu anda ne doğrudur ne de yanlıştır ve “0” ya da
yanlışlık ile “1” ya da doğruluktan farklı üçüncü bir değere sahip olmalıdır.
Bu değeri “$\frac{1}{2}$” ile gösterebiliriz. O, “mümkün olan”ı temsil eder
ve “doğru” ile “yanlış”ın yanına üçüncü bir değer olarak katılır.
\end{quote}
\L ukasiewicz'in üçüncü doğruluk değeri için $\Undef$ kullanacağız.
\footnote{\L ukasiewicz burada “mümkün”ü günümüzde alışılmadık olan bir
anlamda, yani mümkün fakat zorunlu değil anlamında kullanır.}

Bu yorumda $\lnot$, $\land$ ve $\lor$ bağlaçlarının doğruluk fonksiyonlarını
belirlemek kolaydır: geleceğe ilişkin olumsal bir tümcenin değillemesi de
geleceğe ilişkin olumsal bir tümcedir; dolayısıyla
$\tf{\lnot}(\Undef) = \Undef$. Bir tümel evetlemenin bir bileşeni
belirlenmemiş, öteki ise doğruysa tümel evetleme de belirlenmemiştir; sonuçta
geleceğe ilişkin olumsal bileşenin nasıl sonuçlanacağına bağlı olarak
tümel evetleme doğru da yanlış da çıkabilir. Öyleyse
\[
    \tf{\land}(\True, \Undef) =
\tf{\land}(\Undef, \True) =
\Undef.
\]
Öteki bileşen yanlışsa tümel evetleme doğru çıkamaz; dolayısıyla
\[\tf{\land}(\False, \Undef) =
\tf{\land}(\Undef, \False) = \False.\]
Öteki değerler (argümanlar belirlenmiş doğruluk değerleri $\True$ ya da
$\False$ ise) klasik mantıktaki gibidir.

Koşulluda durum biraz daha güçtür. $q$'nun geleceğe ilişkin olumsal bir
önerme olduğunu varsayın. $p$ yanlışsa $q$ nasıl sonuçlanırsa sonuçlansın
$p \lif q$ doğru olacaktır; öyleyse
$\tf{\lif}(\False, \Undef) = \True$ koymalıyız. $p$ doğruysa $q$ ne olursa
olsun $q \lif p$ doğru olacaktır; dolayısıyla
$\tf{\lif}(\Undef, \True) = \True$. $p$ doğruysa $p \lif q$ doğru da yanlış
da çıkabilir; öyleyse $\tf{\lif}(\True, \Undef) = \Undef$. Benzer biçimde
$p$ yanlışsa $q \lif p$ doğru da yanlış da çıkabilir; dolayısıyla
$\tf{\lif}(\Undef, \False) = \Undef$. Geriye hem $p$'nin hem de $q$'nun
geleceğe ilişkin olumsal olduğu durum kalır. Güdülenmeye dayanarak bu durumda
aslında $\Undef$ değeri vermeliyiz. Ancak bu, $!A \lif !A$'yı bir totoloji
\emph{olmaktan çıkarırdı}. \L ukasiewicz, $!A \lor \lnot !A$ ve
$\lnot(!A \land \lnot !A)$'dan vazgeçmekte sakınca görmemiş, fakat
$!A \lif !A$'dan vazgeçmeyi kabul etmemiştir. Bu nedenle
$\tf{\lif}(\Undef, \Undef) = \True$ koşulunu koymuştur.

\begin{defn}\ollabel{def:lukasiewicz}
Üç değerli \L ukasiewicz mantığı şu matrisle tanımlanır:
\begin{enumerate}
  \item $\lnot$, $\land$, $\lor$, $\lif$ bağlaçlarını içeren standart
  önerme dili $\Lang L_0$.
  \item Doğruluk değerleri kümesi $V = \{\True, \Undef, \False\}$.
  \item Tek belirlenmiş değer $\True$'dur; yani $V^+ = \{\True\}$.
  \item Doğruluk fonksiyonları aşağıdaki tablolarla verilir:
  \begin{center}
    \begin{tabular}{c|c} 
      $\tf{\lnot}$ & \\ 
      \hline  
      $\True$ & $\False$ \\ 
      $\Undef$ & $\Undef$ \\
      $\False$ & $\True$ 
    \end{tabular}
    \quad
    \begin{tabular}{c|ccc} 
      $\tf{\land}[\LogLuk[3]]$ & $\True$ & $\Undef$ & $\False$ \\ 
      \hline 
      $\True$ & $\True$ & $\Undef$ & $\False$ \\ 
      $\Undef$ & $\Undef$ & $\Undef$ & $\False$\\ 
      $\False$ & $\False$ & $\False$ & $\False$ 
    \end{tabular}
    \\[2ex]
    \begin{tabular}{c|ccc} 
      $\tf{\lor}[\LogLuk[3]]$ & $\True$ & $\Undef$ & $\False$ \\ 
      \hline 
      $\True$ & $\True$ & $\True$ & $\True$ \\ 
      $\Undef$ & $\True$ & $\Undef$ & $\Undef$ \\
      $\False$ & $\True$ & $\Undef$ & $\False$ 
    \end{tabular}
    \quad
    \begin{tabular}{c|ccc} 
      $\tf{\lif}[\LogLuk[3]]$ & $\True$ & $\Undef$ & $\False$ \\ 
      \hline 
      $\True$ & $\True$ & $\Undef$ & $\False$ \\ 
      $\Undef$ & $\True$ & $\True$ & $\Undef$  \\ 
      $\False$ & $\True$ & $\True$ & $\True$ 
    \end{tabular}
  \end{center} 
\end{enumerate}
\end{defn}

Kolayca görülebileceği gibi yalnızca $\lnot$, $\land$ ve $\lor$ içeren
herhangi bir !!{formula}~$!A$, bütün !!{propositional variable}s{}ne
$\Undef$ atanırsa $\Undef$ doğruluk değerini alır. Dolayısıyla örneğin klasik
totolojiler $p \lor \lnot p$ ve $\lnot(p \land \lnot p)$,
$\pAssign v(p) = \Undef$ olduğunda her zaman
$\pValue{v}(!A) = \Undef$ olduğu için $\LogLuk[3]$ içinde totoloji değildir.

$\pAssign v(p) = \True$ ya da $\False$ olan !!{valuation}s{}de
$\pValue v(!A)$, klasik doğruluk değeriyle çakışır.

\begin{prop}
  $!A$ içindeki bütün $p$'ler için
  $\pAssign v(p) \in \{\True, \False\}$ ise
  $\pValue v(!A)[\LogLuk[3]] = \pValue v(!A)[\LogCL]$.
\end{prop}

\begin{prob}\label{mvl:thr:luk:prob:luk-iff} $\LogLuk[3]$ içinde
  $\pValue v(!A \liff !B) = \pValue v((!A \lif !B) \land (!B \lif
  !A))$ olarak tanımladığımızı varsayın. $\liff$ hangi doğruluk tablosuna
  sahip olur?
\end{prob}

Birçok klasik totoloji, örneğin $\lnot p \lif (p \lif q)$, aynı zamanda
\LogLuk[3] içinde de totolojidir. Klasik mantıkta olduğu gibi bunu doğruluk
tablolarıyla doğrulayabiliriz:
\begin{center}
\begin{tabular}{cc|cccccc}
  $p$ & $q$ & $\lnot$ & $p$ & $\lif$ & $\smash{(}p$ & $\lif$ & $q\smash{)}$ \\ \hline
  \True & \True & \False & \True & \True & \True & \True & \True \\
  \True & \Undef & \False & \True & \True & \True & \Undef & \Undef \\
  \True & \False & \False & \True & \True & \True & \False & \False \\
  \Undef & \True & \Undef & \Undef & \True & \Undef & \True & \True \\
  \Undef & \Undef & \Undef & \Undef & \True & \Undef & \True & \Undef \\
  \Undef & \False & \Undef & \Undef & \True & \Undef & \Undef & \False \\
  \False & \True & \True & \False & \True & \False & \True & \True \\
  \False & \Undef & \True & \False & \True & \False & \True & \Undef \\
  \False & \False & \True & \False & \True & \False & \True & \False \\  
\end{tabular}
\end{center}

\begin{prob}
  Aşağıdakilerin \LogLuk[3] içinde totoloji olduğunu gösterin:
  \begin{enumerate}
    \item $p \lif (q \lif p)$
    \item\label{mvl:thr:luk:prob:luk-taut-2} $\lnot(p \land q) \liff (\lnot p \lor \lnot q)$
    \item\label{mvl:thr:luk:prob:luk-taut-3} $\lnot(p \lor q) \liff (\lnot p \land \lnot q)$
  \end{enumerate}
  (\olref[mvl][thr][luk]{prob:luk-taut-2} ve
  \olref[mvl][thr][luk]{prob:luk-taut-3} içinde $!A \liff !B$'yi
  $(!A \lif !B) \land (!B \lif !A)$'nın bir kısaltması olarak alın ya da
  \olref[mvl][thr][luk]{prob:luk-iff} için verdiğiniz çözüme başvurun.)
\end{prob}

\begin{prob}
  Aşağıdaki klasik totolojilerin \LogLuk[3] içinde totoloji olmadığını
  gösterin:
  \begin{enumerate}
    \item $(\lnot p \land p) \lif q$
    \item $((p \lif q) \lif p) \lif p$
    \item $(p \lif (p \lif q)) \lif (p \lif q)$
  \end{enumerate}
\end{prob}

Bu nedenle bütün klasik totolojiler $\LogLuk[3]$'te totoloji olmasa bile her
!!{valuation} altında en azından ya $\True$ ya da $\Undef$ değerini almaları
gerektiği düşünülebilir. Oysa durum böyle değildir. Bunun bir karşı örneği
\[
  \lnot(p \lif \lnot p) \lor \lnot(\lnot p \lif p)
\]
formülüdür; bu formül $p$, $\Undef$ olduğunda $\False$ olur.

\begin{prob}
  Aşağıdaki bağıntılardan hangileri \L ukasiewicz mantığında geçerlidir? Her
  biri için bir doğruluk tablosu verin.
  \begin{enumerate}
    \item $p, p \lif q \Entails q$
    \item $\lnot\lnot p \Entails p$
    \item $p \land q \Entails p$
    \item $p \Entails p \land p$
    \item $p \Entails p \lor q$
  \end{enumerate}
\end{prob}

\L ukasiewicz, bir-konumlu $\Diamond !A$ (“$!A$ mümkündür”) bağlacıyla ona
karşılık gelen $\Box !A$ (“$!A$ zorunludur”) bağlacını ekleyerek üç değerli
sistemine dayalı bir olanak mantığı kurmayı umuyordu:
\begin{center}
  \begin{tabular}{c|c} 
    $\tf{\Diamond}$ & \\ 
    \hline  
    $\True$ & $\True$ \\ 
    $\Undef$ & $\True$ \\
    $\False$ & $\False$ 
  \end{tabular}
  \quad
  \begin{tabular}{c|c} 
    $\tf{\Box}$ & \\ 
    \hline  
    $\True$ & $\True$ \\ 
    $\Undef$ & $\False$ \\
    $\False$ & $\False$ 
  \end{tabular}
\end{center}
Başka bir deyişle $p$ ancak ve ancak yanlış olduğu henüz belirlenmemişse
mümkündür; ancak ve ancak doğru olduğu zaten belirlenmişse zorunludur.

\begin{prob}
  $\Box p \liff \lnot\Diamond \lnot p$ ile $\Diamond p \liff
  \lnot \Box \lnot p$'nin, $\Box$ ve $\Diamond$ doğruluk tablolarıyla
  genişletilmiş \LogLuk[3] içinde totoloji olduğunu gösterin.
\end{prob}

Ancak önerilen bu kiplik mantığının eksiklikleri kısa sürede ortaya çıktı:
Koşullar nasıl sonuçlanırsa sonuçlansın $p \land \lnot p$ hiçbir zaman doğru
olamaz. Dolayısıyla şimdi belirlenmemiş (ve bu nedenle değeri de
belirlenmemiş) olsa bile olanaksız sayılmalıdır; yani
$\lnot \Diamond(p \land \lnot p)$ bir totoloji olmalıdır. Oysa
$\pAssign v(p) = \Undef$ ise
$\pValue v(\lnot \Diamond(p \land \lnot p)) = \False$ olur.
\L ukasiewicz, olanaklılık ve zorunluluk gibi kiplik ayrımlarını barındırmak
için iki doğruluk değerinin yetmeyeceği konusunda haklı olsa da üçüncü bir
doğruluk değeri getirmek de yeterli değildir.

\end{document}
